%%%%%%%%%%%%%%%%%%%%%%%%%%%%%%%%%%%%%%%%
\chapter{Objetivos}
%%%%%%%%%%%%%%%%%%%%%%%%%%%%%%%%%%%%%%%%

En esta sección se presentan los objetivos que orientan el desarrollo de la revisión sistemática. Estos permiten delimitar el alcance del estudio y organizar el análisis de la literatura científica.


%%%%%%%%%%%%%%%%%%%%%%%%%%%%%%%%%%%%%%%%
\section{Objetivo general}

\begin{quote}
    Caracterizar el estado del arte en el reconocimiento de imágenes de polinizadores en aprendizaje profundo, mediante una revisión sistemática de estudios científicos desde 2020 sobre redes neuronales en clasificación, detección o segmentación.
\end{quote}


%%%%%%%%%%%%%%%%%%%%%%%%%%%%%%%%%%%%%%%%
\section{Objetivos específicos}

A partir del objetivo general, se establecen los siguientes objetivos específicos, los cuales permiten descomponer el análisis en dimensiones concretas relacionadas.

\begin{itemize}
  \item Identificar los modelos y arquitecturas de aprendizaje profundo en el reconocimiento de imágenes de polinizadores y las configuraciones de redes neuronales.
  \item Comparar las métricas de desempeño de la literatura científica, tales como precisión, exactitud, sensibilidad y medida F1, de acuerdo con la tarea y el tipo de modelo.
  \item Analizar los factores técnicos que se asocian con el rendimiento de los modelos, como la transferencia de aprendizaje, el aumento de datos y los métodos de optimización.
\end{itemize}